\documentclass[11pt]{article}
\usepackage{mathunal-base}
\mathunalfuente{serif}

\renewcommand{\examtitulo}{Primer Examen Parcial de Álgebra Lineal}
\renewcommand{\examfecha}{24 de Marzo de 2018}

\newcommand{\vf}{\fbox{V}\ \fbox{F}}

\begin{document}

\examheader

\vspace{4pt}
\noindent
\begin{minipage}[t]{0.62\linewidth}
\textbf{Puntaje. Sólo para uso Oficial}\\[4pt]
\begin{tabular}{|c|c|c|c|c|c|c|}
\hline
1--5 & 6 & 7 & 8 & 9 & \textbf{TOTAL} & \textbf{NOTA} \\ \hline
 & & & & & & \\ \hline
\end{tabular}
\end{minipage}

\vspace{6pt}
\noindent\textbf{Instrucciones:} La duración del examen es de 1 hora y 50 minutos. El examen consta de nueve preguntas en tres hojas impresas por ambos lados, verifique que su examen esté completo. En las preguntas con procedimiento justifique sus respuestas en los espacios asignados. No está permitido sacar hojas en blanco ni ningún tipo de apuntes durante el examen, verifique que su celular esté apagado. No se permite el uso de calculadora.

\vspace{6pt}
\noindent\textbf{IDENTIFICACIÓN}

\vspace{8pt}
\noindent Nombre: \dotfill\ Cédula \dotfill

\vspace{4pt}
\noindent Profesor: \dotfill\ Grupo \dotfill

\vspace{6pt}
\begin{center}
\textbf{Espacio en blanco para realizar cómputos}
\end{center}

\newpage
\noindent\textbf{\large I. Completación}

\vspace{2pt}
\noindent En las preguntas 1 a 5 complete los espacios en blanco. \textbf{NOTA}: En esta sección se califica sólo la respuesta y no se tiene en cuenta el procedimiento.

\begin{enumerate}
\item{}[6pt] Sean $\mathbf{u}, \mathbf{v}$ vectores unitarios en $\mathbb{R}^n$ tales que $\mathbf{u}\cdot(2\mathbf{v})=\sqrt{3}$. El ángulo entre $\mathbf{u}$ y $\mathbf{v}$ es

\vspace{4pt}
\fbox{\begin{minipage}[t][2.3cm]{\dimexpr\linewidth-2\fboxsep-2\fboxrule}
$\theta = $
\end{minipage}}

\item{} Sean $A$ y $B$ matrices $n\times n$. Señale el valor de verdad de cada enunciado.
\begin{enumerate}
\item{}[2pt] Si $AB=BA$, entonces $(A-B)(A+B)=A^2-B^2$. \hfill\vf
\item{}[2pt] Si $A$ y $B$ son simétricas, entonces $A^3-5B$ también es simétrica. \hfill\vf
\item{}[2pt] Si $AB$ es invertible, entonces $A$ y $B$ también son invertibles. \hfill\vf
\item{}[2pt] El determinante de $A+I_n$ es igual a $\det A + 1$. \hfill\vf
\end{enumerate}

\item{} Suponga que $\begin{vmatrix}a&b&c\\d&e&f\\g&h&i\end{vmatrix}=5$. Calcule los siguientes determinantes

\vspace{6pt}
\noindent
\begin{minipage}[t]{0.46\linewidth}
(a) [2pt] $\begin{vmatrix}c&b&a\\f&e&d\\i&h&g\end{vmatrix} = $ \dotfill
\end{minipage}%
\hfill
\begin{minipage}[t]{0.5\linewidth}
(b) [2pt] $\begin{vmatrix}a&b&c\\d+g&e+h&f+i\\g&h&i\end{vmatrix} = $ \dotfill
\end{minipage}

\vspace{10pt}
\noindent (c) [4pt] $\begin{vmatrix}4b&4e&4h\\a&d&g\\c&f&i\end{vmatrix} = $ \dotfill

\item{}[8pt] Sean $v_1 = \begin{bmatrix}1\\0\\1\end{bmatrix}$, $v_2 = \begin{bmatrix}0\\1\\2\end{bmatrix}$ y $v_3 = \begin{bmatrix}1\\1\\k\end{bmatrix}$. Halle todos los valores de $k$ tales que $v_3\in gen(v_1,v_2)$.

\vspace{4pt}
\fbox{\begin{minipage}[t][2.3cm]{\dimexpr\linewidth-2\fboxsep-2\fboxrule}
$k = $
\end{minipage}}

\item{}[10pt] Halle un conjunto $\mathcal{B}$ en $\mathbb{R}^4$ tal que $gen(\mathcal{B})$ sea igual al conjunto solución del sistema lineal homogéneo
\[
\left\{
\begin{aligned}
x+2y+2z-w &= 0\\
3x+7y-6z+5w &= 0
\end{aligned}
\right.
\]

\vspace{4pt}
\fbox{\begin{minipage}[t][3.5cm]{\dimexpr\linewidth-2\fboxsep-2\fboxrule}
$\mathcal{B} = $
\end{minipage}}
\end{enumerate}

\vspace{6pt}
\begin{center}\textbf{II. Solución con Procedimiento}\end{center}

\begin{enumerate}
\setcounter{enumi}{5}
\item{}[10pt] Sea $\{\mathbf{u},\mathbf{v},\mathbf{w}\}$ un conjunto linealmente independiente de vectores en $\mathbb{R}^n$. Demostrar que el conjunto $\{\mathbf{u},\mathbf{v},2\mathbf{u}+\mathbf{v}-3\mathbf{w}\}$ también es linealmente independiente.

\vspace{6cm}

\item{}[25pt] Un departamento de pesca y caza del estado proporciona tres tipos de comida a un lago que alberga a tres especies de peces. Cada pez de la especie 1 consume cada semana un promedio de 1 unidad del alimento A, 1 unidad del alimento B y 2 unidades del alimento C. Cada pez de la especie 2 consume cada semana un promedio de 3 unidades del alimento A, 4 del B y 5 del C. Para la especie 3, el promedio semanal de consumo por pez es de 2 unidades del alimento A, 1 unidad del B y 5 del C. Cada semana se proporcionan al lago 25\,000 unidades del alimento A, 20\,000 unidades del B y 55\,000 del C.
\begin{enumerate}
\item{}[17pt] Si suponemos que los peces se comen todo el alimento, plantee y resuelva un sistema de ecuaciones lineales que permita hallar las cantidades de cada especie que pueden coexistir en el lago.

\vspace{6cm}

\item{}[8pt] Encuentre los intervalos de variación de cada variable, es decir, encuentre los valores máximos y mínimos de peces de cada especie que puede haber.

\vspace{4cm}
\end{enumerate}

\item{}[10pt] Encontrar el valor máximo de la función $f(x,y)=2x+3y$ sujeto a las restricciones $0\leq x\leq 30$, $0\leq y\leq 20$ y $x+2y\leq 54$.

\vspace{6cm}

\item{}[15pt] Considere la matriz invertible $A = \begin{bmatrix}2&1&3\\-2&1&-6\\1&-1&4\end{bmatrix}$.

Use el método de Gauss-Jordan para hallar la inversa de $B=\dfrac{1}{500}A^T$.
\end{enumerate}

\examfooter
\end{document}
